

\begin{notebox}
This material was written while I was first working with stickelbacks and before I realised
I wanted to turn them about and re-index them from one not zero. If any of this material
is used in future it needs to be re-jigged 
\end{notebox}
\subsection{Tightening}
\workt{Having defined the initial composition of tight zigzags:}
EXCEPT that the above composition is not necessarily tight because
it will not necessarily be the case that $\overline{f_{n-1}}=\overline{f_n\circ g_1}$.

Therefore we have to do some modifications to both halves of the composition.

\begin{oldtt}
Modifications to first (f,f') half will be along these lines:

$$
\arraycolsep=7pt
\begin{array}{c c c c c}
\Rnode{xnpp}{x_{n-2}} && \kern0.25cm\Rnode{xnp}{x_{n-1}} && \kern0.5cm\Ob{z}{0}\\[1cm]
  & \kern0.25cm\Rnode{wnp}{w_{n-1}}  && \kern0.5cm\Ob{w}{n}
\end{array}
\begin{arrows}
\ncarr{z0}{wn}\alabel{g_1 f_n}[0.25][0.4]
\ncarr{xnp}{wn}\alabel{f'_{n-1}.\reallywidehat{g_1 f_n}}[0.4][0.4]
\ncarr{xnp}{wnp}\alabel{f_n.\overline{f'_{n-1}.\reallywidehat{g_1 f_n}}}[0.5][0.4]
\ncarr{xnpp}{wnp}\blabel{f'_{n-2}.\reallywidehat{f_n.\overline{f'_{n-1}\reallywidehat{g_1 f_n}}}}[0.45][0.4]
\end{arrows}
$$
\end{oldtt}
\begin{oldtt}
\textbf{TIGHTENING a loose stickleback.}

start with

$$
\arraycolsep=7pt
\begin{array}{cccccc}
 &\bOb{v}& &\bOb{x}& &\bOb{z} \\[0.8cm]
\bOb{u}& &\bOb{w}& &\bOb{y}&
\end{array}
\begin{arrows}
\ncarr{v}{u}\blabel{e}[0.25][0.4]
\ncarr{v}{w}\alabel{f}[0.25][0.4]
\ncarr{x}{w}\blabel{g}[0.25][0.4]
\ncarr{x}{y}\alabel{h}[0.25][0.4]
\ncarr{z}{y}\blabel{i}[0.25][0.4]
\end{arrows}
$$

This tightens to 

$$
\arraycolsep=7pt
\begin{array}{cccccc}
 &\bOb{v}& &\bOb{x}& &\bOb{z} \\[0.8cm]
\bOb{u}& &\bOb{w}& &\bOb{y}&
\end{array}
\begin{arrows}
\ncarr{v}{u}\blabel{r_e}[0.25][0.4]
\ncarr{v}{w}\alabel{r_f}[0.25][0.4]
\ncarr{x}{w}\blabel{r_g}[0.25][0.4]
\ncarr{x}{y}\alabel{r_h}[0.25][0.4]
\ncarr{z}{y}\blabel{r_i}[0.25][0.4]
\end{arrows}
$$


where
\newcommand{\fConstriction}
{\reallywidehat{\overline{h\circ\widehat{i}}\circ g}}
\newcommand{\eRestriction}
{\overline{f\circ\fConstriction}}
\newcommand{\rebar}{\overline{e} \circ \eRestriction}
\newcommand{\rf}{\overline{e} \circ f \circ \reallywidehat{\overline{h \circ \widehat{i}} \circ g} }
\newcommand{\rfhat}{\widehat{\overline{e} \circ f} \circ \fConstriction}
\newcommand{\hRestriction}
{\overline{g \circ \reallywidehat{\overline{e} \circ f}}}
\newcommand{\rgbar}{\hRestriction\circ\overline{h\circ\widehat{i}}}
\newcommand{\rhhat}{\widehat{i} \circ \reallywidehat{\hRestriction \circ h}}
\newcommand{\ri}
{i \circ \reallywidehat{\hRestriction\circ h}}
\newcommand{\hRestrictionSimplified}
{\overline{g \circ \widehat{f}}}
\begin{align*}
r_e&= \eRestriction \circ e  
& \overline{r_e} &= \overline{e} \circ \eRestriction 
&\widehat{r_e}&=\reallywidehat{\eRestriction \circ e}\\
r_f&=\rf 
& \overline{r_f} &= \rebar
&\widehat{r_f} &= \rfhat\\
r_g&=\overline{h \circ \widehat{i}} \circ g \circ \reallywidehat{\overline{e} \circ f}
& \overline{r_g} &= \rgbar
&\widehat{r_g} &= \rfhat   \\
r_h&=\hRestriction\circ h\circ\widehat{i}
&\overline{r_h}&=\rgbar
&\widehat{r_h}&=\rhhat\\
r_i&=\ri&\overline{r_i}&=\overline{\ri}
&\widehat{r_i}&=\rhhat
\end{align*}
Note that $\overline{r_e}, \overline{r_f}, \overline{r_g}$ and $\overline{r_h}$ have been variously simplified using R.2 and R.3 and that $\overline{r_i}$ does not simplify. Likewise $\widehat{r_f}, \widehat{r_g} , \widehat{r_h}$ and $\widehat{r_i}$ have been variously simplified using lemma \lref{fglemma}, (i) and (iii), and $\widehat{r_e}$ does not simplify.

\textbf{WHAT HAPPENS IF ONLY THE LAST LEG FAILS THE TIGHTNESS CONDITION} 
Suppose $\overline{e}=\overline{f}$, $\widehat{f}=\widehat{g}$, $\overline{g}=\overline{h}$. Then we get the some simplifications as follows.

Stage one $bar{e}=\overline{f}$: 
\begin{align*}
&&& \overline{e}=\overline{f} && \widehat{f} = \widehat{g}  && \overline{g} = \overline{h} \\
r_e&= \eRestriction \circ e  
&& 
&& \\
r_f&=\rf  
&&= f \circ \reallywidehat{\overline{h \circ \widehat{i}} \circ g} 
&& \\
r_g &=\overline{h \circ \widehat{i}} \circ g \circ \reallywidehat{\overline{e} \circ f}
&&= \overline{h \circ \widehat{i}} \circ g \circ \reallywidehat{f}
&&=  \overline{h \circ \widehat{i}} \circ g \\
r_h &=\hRestriction\circ h\circ\widehat{i}
&&= \hRestrictionSimplified\circ h\circ\widehat{i} 
&&= \overline{g} \circ h \circ \widehat{i}
&&= h \circ \widehat{i} \\
r_i &=\ri 
&&=i \circ \reallywidehat{\hRestrictionSimplified \circ h}
&&=i \circ \reallywidehat{\overline{g} \circ h}
&&= i \circ \widehat{h}
\end{align*}

\textbf{Rename}
Rename $e,f,g,h,i$ to be $f_0,f'_0,f_1,f'_1,f_2$.

Rename $r_e,r_f,r_g,r_h,r_i$ to be $t_0,t'_0,t_1,t'_1,t_2$.

Then  \textbf{what happens if only the last leg fails the tightness condition:}

\begin{alignat*}{2}
t_2 &=f_2 \circ \widehat{f'_1} \\
t'_1&=f'_1\circ\widehat{t_2} &&= f'_1 \circ \widehat{f_2}\\
t_1&= \overline{t'_1}\circ f_1&=\overline{f'_1 \circ \widehat{f_2}} \circ f_1\\
t'_0&=f'_0 \circ \widehat{t_1}=f'_0 \circ \reallywidehat{\overline{f'_1 \circ \widehat{f_2}} \circ f_1}\\
t_0&=\overline{t'_0}\circ f_0
\end{alignat*}

\textbf{Same but with just $f_0,f'_0$ and $f_1$}

\begin{alignat*}{2}
t_1&= f_1 \circ \widehat{f'_0}\\
t'_0&=f'_0 \circ \widehat{t_1} &&=f'_0 \circ \widehat{f_1}\\
t_0&=\overline{t'_0} \circ f_0 &&= \overline{f'_0 \circ \widehat{f_1}} \circ f_0 
\end{alignat*}

\textbf{Same but substitute $g_0 \circ f_1$ for $f_1$.}
\begin{alignat*}{2}
t_1&= g_0 \circ f_1 \circ \widehat{f'_0}  &&= g_0 \circ f_1\\
t'_0&=f'_0 \circ \widehat{t_1} &&=f'_0 \circ \widehat{g_0 \circ f_1}\\
t_0&=\overline{t'_0} \circ f_0 &&= \overline{f'_0 \circ \widehat{g_0 \circ f_1}} \circ f_0 
\end{alignat*}
\end{oldtt}

\subsection{Trying and failing to prove R.4 }
\subsubsection{Before the change of direction of depiction}
to prove R.4 we need first prove Cockett et als.: wR.4 which is that

\begin{equation}
\label{restrictionlemmiiasR4Antecedent}
\overline{F \circ G}=\overline{F \circ \overline{G}}
\end{equation}


we can then proceed as follows.

\textbf{init(F) :: last(F)}

First consider that if $F$ is the zigzag
$$
\arraycolsep=7pt
\zigzag{x}{y}{f}{n}
$$
then if $n \geq 2$ then $F$ 
then it can be expressed as  the composition of zigzags
\begin{equation}
\label{FheadTail}
F=init(F)::last(F)
\end{equation}
where
\begin{align*}
init(F) &= \tuple{f_0,f'_0,...,f'_{n-2},f_{n-1}} \\
last(F) &= \tuple {\overline{f'_{n-1}},f'_{n-1},f_n}
\end{align*}

We will first prove that R.4 holds for $F$ of length $2$, the case of $n=2$,
and then prove R.4 holds for $F$ of length greater than $2$ by induction on $n$.

\textbf{Case $n=2$}

Suppose $that F$ is the zigzag
$$
\arraycolsep=7pt
\begin{array}{cccc}
 &\bOb{x}& &\bOb{z} \\[0.8cm]
\bOb{w}& &\bOb{y}&
\end{array}
\begin{arrows}
\ncarr{x}{w}\blabel{f_0}[0.25][0.4]
\ncarr{x}{y}\alabel{f'_0}[0.25][0.4]
\ncarr{z}{y}\blabel{f_1}[0.25][0.4]
\end{arrows}
$$
of length $2$.
We need show that

$$
F \circ \overline{G}  = \overline{F \circ G} \circ F 
$$


\begin{align*}
LHS &= \tuple{f_0,f'_0,f_1} \circ \overline{G} \\
    &= T(\tuple{f_0,f'_0,\reallywidehat{g_0} \circ f_1}) &&\mbox{using the definition of the $\bar{}$ operator for zigzags} \\
    &= \tuple{ \overline {f'_0  \circ \reallywidehat{g_0 \circ f_1}} \circ f_0, 
                 f'_0  \circ \reallywidehat{g_0 \circ f_1}, 
                 \reallywidehat{g_0} \circ f_1   
                }  && \mbox{by earlier that needs be lemmarised, and by RR.4 (twice)}
\end{align*}

\begin{align*}
RHS &=  \overline{F \circ G} \circ F \\
    &= \overline{F \circ \overline{G}} \circ F 
                  &&\mbox{ using (\ref{restrictionlemmiiasR4Antecedent}),}\\
    &= \overline{\tuple{ \overline {f'_0  \circ \reallywidehat{g_0 \circ f_1}} \circ f_0, 
                 f'_0  \circ \reallywidehat{g_0 \circ f_1}, 
                 \reallywidehat{g_0} \circ f_1   
                }} \circ F 
                  &&\mbox{using LHS as above,} \\
    &= \reallywidehat{\overline {f'_0  \circ \reallywidehat{g_0 \circ f_1}} \circ f_0} \circ F
                  && \mbox{ from definition of the $\bar{\ }$ operator,} \\
    &= T(\tuple{
            f_0 \circ\reallywidehat{\overline {f'_0  \circ \reallywidehat{g_0 \circ f_1}} \circ f_0},
            f'_0,
            f_1
        })         
                  && \mbox{from definition of $\circ$}\\
    &= \tuple{
            t_0,
            \overline{t_0} \circ f'_0,
            f_1 \circ \reallywidehat{\overline{t_0} \circ f'_0}
        } \mbox{ where } t_0 = f_0 \circ\reallywidehat{\overline {f'_0  \circ \reallywidehat{g_0 \circ f_1}} \circ f_0} 
    && \mbox{ by lemma \lref{leftLowLemma} leftLowLemma}\\
    &= &&\mbox{something that  hopefully matches rhs of LHS, above}   \\
    &= \mbox{\highlight{but it doesn't equate to LHS and we fail!}}
\end{align*}

\textbf{case $n > 2$}
For $F$ of length greater than $2$ we prove by induction. The inductive hypothesis is
\begin{equation}
\label{R4InductiveHypothesis}
init(F) \circ \overline{G} = \overline{init(F) \circ G} \circ init(F)
\end{equation}
Note first that
\begin{equation}
\label{headTailAndG} 
init(F) \circ tail(F) \circ G= F \circ G
\end{equation} 
follows from the definitions and from lemma \lref{tighteningCommutesWithCompositionLemma}.

\begin{align*}
F \circ \overline{G} 
    &= (init(F) \conc last(F)) \circ \overline{G}
                            && \mbox{using (\ref{FheadTail}),}\\
    &= T( init(F) \conc last(F) \conc \overline{G}) 
                            && \mbox{from definition of $\circ$ and assicativity of $\conc$,}\\
    &= T( init(F) \conc T(last(F) \conc \overline{G}) )
                            && \mbox{by lemma \lref{tighteningCommutesWithCompositionLemma},}\\
    &= T( init(F) \conc (last(F) \circ \overline{G}) )
                            && \mbox{using definition of $\circ$,}\\
    &= T(init(F) \conc (\overline{last(F) \circ G} \circ last(F)) )
                            && \mbox{using $n=1$ case,} \\
    &= init(F) \circ \overline{last(F) \circ G} \circ last(F) 
                            && \mbox{using definition of $\circ$,} \\
    &= \overline{init(F) \circ last(F) \circ G} \circ init(F)  \circ last{F} 
                            && \mbox{using the inductive hypothesis (\ref{R4InductiveHypothesis}),} \\
    &= \overline{F \circ G} \circ F 
                                && \mbox{as required, by applying (\ref{headTailAndG}).}
\end{align*}
\subsubsection{Turn zigzags around}
Looking at Cockett's Shein's lemma construction make me think I have zigzags the wrong way around.
That hios zigzag from x to y should be a morphism from y to x in the category of zigzags.
I can't see this making a difference but I need to check out the preceding failed lemma to see if it succeeds when things are turned around. I can't see how this can be but I better try it anyway. I will edit the previous by reversing the directions of the arrows. It is better this way anyway because the embedding from \catcw into zigzags of \catcw is a more obvious. 
Start with the bit  that I couldn't prove and see if I can now prove it even though I don't see how it can have changed significantly so that it goes through.
\textbf{prove that R.4 holds for $F$ of length 2.}
Suppose $that F$ is the zigzag
$$
\arraycolsep=7pt
\begin{array}{cccc}
 &\bOb{x}& &\bOb{z} \\[0.8cm]
\bOb{w}& &\bOb{y}&
\end{array}
\begin{arrows}
\ncarr{w}{x}\alabel{f_0}[0.5][0.4]
\ncarr{y}{x}\blabel{f'_0}[0.5][0.4]
\ncarr{y}{z}\blabel{f_1}[0.5][0.4]
\end{arrows}
$$
of length $2$.
We need show that

$$
F \circ \overline{G}  = \overline{F \circ G} \circ F 
$$


\highlight{the following looks right}
\begin{align*}
LHS &= \tuple{f_0,f'_0,f_1} \circ \overline{G} \\
    &= T(\tuple{f_0,f'_0, f_1 \circ \overline{g_0}}) 
                   &&\mbox{using the definition of the $\bar{}$ operator for zigzags} \\
    &= \tuple{f_0 \circ \reallywidehat {\overline{f_1 \circ g_0} \circ f'_0}, \ 
                 \overline{f_1 \circ g_0 } \circ f'_0, \
                 f_1 \circ \overline{g_0}   
                }  && \mbox{introduce some earlier lemma to cover this}
\end{align*}

\highlight{begin editing }
\begin{align*}
RHS &=  \overline{F \circ G} \circ F \\
    &= \overline{F \circ \overline{G}} \circ F 
                  &&\mbox{ using a (\ref{restrictionlemmiiasR4Antecedent}) CHECK --- THIS IS wR4,}\\
    &= \overline{\tuple{f_0 \circ \reallywidehat {\overline{f_1 \circ g_0} \circ f'_0}, \ 
                 \overline{f_1 \circ g_0 } \circ f'_0, \
                 f_1 \circ \overline{g_0}   
                }} \circ F 
                  &&\mbox{using LHS as above,} \\
    &= \tuple{\overline{f_0 \circ \reallywidehat {\overline{f_1 \circ g_0} \circ f'_0}}} \circ F
                  && \mbox{ from definition of the $\bar{\ }$ operator,} \\
    &= T(\tuple{
            \overline{f_0 \circ \reallywidehat {\overline{f_1 \circ g_0} \circ f'_0}}
                                \circ f_0,\ 
            f'_0,\ 
            f_1
        })         
                  && \mbox{from definition of $\circ$}\\
    &= \tuple{
            t_0,
            f'_0 \circ \widehat{t_0} ,
            f_1 \circ \overline{\widehat{f'_0 \circ t_0}}
        } \mbox{ where } 
        t_0 = \overline{f_0 \circ \reallywidehat {\overline{f_1 \circ g_0} \circ f'_0}} \circ f_0 
    && \mbox{ by lemma \lref{leftLowLemma} leftLowLemma OR SUCH} \\
    &= &&\mbox{something that  hopefully matches rhs of LHS, above}   \\
    &= \mbox{\highlight{does it equate to LHS or we fail!}}
\end{align*}

We can prove using R.3 and RR.1, for starters, that
\begin{equation}
\overline{f_0 \circ \reallywidehat {\overline{f_1 \circ g_0} \circ f'_0}} \circ f_0 
= f_0 \circ \reallywidehat {\overline{f_1 \circ g_0} \circ f'_0}
\end{equation}

Now the second term of RHS is
\[
f'_0 \circ \reallywidehat{f_0 \circ \reallywidehat {\overline{f_1 \circ g_0} \circ f'_0}}
\]
we need to show this is equal to 
\[
\overline{f_1 \circ g_0 } \circ f'_0
\]

The second term simplifies using fglemma (iii) to
\[
f'_0 \circ \widehat{f_0} \circ \reallywidehat {\overline{f_1 \circ g_0} \circ f'_0}
\]

then simplies to
\[
f'_0 \circ \reallywidehat {\overline{f_1 \circ g_0} \circ f'_0}
\]

which doesn't go thru \highlight{FAILS AS BEFORE} see What isn't true.

\subsection{What Isn't True}

When I tried to prove some property of the stickleback category
(which one!) corresponding to a range category \catcw 
then I needed to 
 prove that in any range category
 $$\overline{g} \circ f \leq f \circ \reallywidehat{\overline{g} \circ f}.$$
This isn't generally true though. We cannot prove it.
\begin{worktt}The fact that it isn't true may become significant.
We may need the fact that this identity doesn't hold in order to show that stickleback category satisfies the awR.4 condition. 
\end{worktt}

This identity doesn't hold 
in the range category of partial sets and functions for suppose
$a$ is a doubleton set ${x_1,x_2}$ 
and that $f$ is defined at both $x_1$ and $x_2$ and that $f(x_1) = f(x_2)=y$.
Suppose $g$ is defined at $x_1$ but not at $x2$. Then
\begin{align*}
x_1 &\xmapsto{f} y \\
x_2 &\mapsto y \\
             \\
x_1 &\xmapsto{\overline{g}\kern0.05cm} x_1 \\
x_2 &\xmapsto{\kern0.2cm} \bot \\
                \\
x_1 &\xmapsto{\overline{g} \circ f} y \\
x_2 &\xmapsto{\kern0.4cm} \bot
\end{align*}

so that $\reallywidehat{\overline{g} \circ f}=\set{y}$  
and

\begin{align*}
x_1 &\xmapsto{f \circ \reallywidehat{\overline{g} \circ f}} y \\
x_2 &\xmapsto {\kern.75cm}y
\end{align*}

and we see that $$\overline{g} \circ f \neq f \circ \reallywidehat{\overline{g} \circ f}.$$
